\section{Introduction}\label{sec:introduction}

The Mahler conjecture asks which convex bodies have the smallest product
of their volume and the volume of a suitably centered polar body.
We resolve the conjecture positively: simplices are exactly the minimizers
in every positive dimension.

A \emph{convex body} in $\mathbb R^n$ is a compact convex set with nonempty
interior.  We write $|A|$ for the $n$-dimensional Lebesgue volume of a measurable
set $A$, and $\langle\cdot,\cdot\rangle$ for the Euclidean inner product.
For $z\in\operatorname{int}K$, define
\[
 (K-z)^\circ
 =\{y\in\mathbb R^n:\langle y,x-z\rangle\le1\text{ for every }x\in K\}.
\]
The minimum of $|(K-z)^\circ|$ is attained at a unique point
$s(K)\in\operatorname{int}K$, the \emph{Santal\'o point}
(a proof is included in Section~\ref{sec:01-cones}). The volume product is
\[
 P(K)=|K|\,|(K-s(K))^\circ|
     =\inf_{z\in\operatorname{int}K}|K|\,|(K-z)^\circ|.
\]
It is invariant under invertible affine maps.
A simplex is the convex hull of $n+1$ affinely independent points.

\begin{theorem}[General Mahler conjecture]\label{thm:mahler}
For every integer $n\ge1$ and every convex body $K\subset\mathbb R^n$,
\[
                 P(K)\ge\frac{(n+1)^{n+1}}{(n!)^2}.
\]
Equality holds if and only if $K$ is a simplex.
\end{theorem}

Theorem~\ref{thm:mahler} applies without symmetry or boundary regularity assumptions.
It also gives the displayed lower bound for
$|K|\,|(K-z)^\circ|$ at every interior translation point $z$.
The symmetric Mahler conjecture asks for the different, stronger constant
$4^n/n!$ on the restricted class of centrally symmetric bodies; that sharper
statement is separate from Theorem~\ref{thm:mahler} and is proved in the
companion paper \cite[Theorem~1.1]{OpenAISymmetricMahler2026}.

\subsection*{History and context}

The lower-volume-product problem grew out of Mahler's work on polarity
and the geometry of numbers. His polygon paper established the general
planar inequality and its triangle equality case among polygons
\cite{Mahler1938}; his transference paper formulated the
higher-dimensional symmetric problem \cite{Mahler1939}.
Meyer later gave a new proof of the planar general inequality and
completed its equality characterization for arbitrary planar convex
bodies: triangles are exactly the minimizers \cite{Meyer1991}. The affine center used here belongs to the classical
Santal\'o theory \cite{Santalo1949}; we give its elementary existence
and uniqueness argument in Section~\ref{sec:01-cones}.

Several lines of work explain the distinction between an asymptotic lower
bound and the exact simplex constant. The inverse Santal\'o theorem of
Bourgain and Milman gives
$|K|\,|K^\circ|\ge c^n|B_2^n|^2$ for origin-symmetric bodies, where $c>0$ is
absolute and $B_2^n$ is the Euclidean unit ball
\cite{BourgainMilman1987}. Kuperberg obtained explicit lower bounds through
Gauss linking integrals, with separate conclusions in the symmetric and
general settings \cite{Kuperberg2008}. Nazarov developed a complex-analytic
proof using Bergman kernels and H\"ormander estimates
\cite{Nazarov2012}; Mastrantonis and Rubinstein extended that direct
analytical method to nonsymmetric bodies \cite{MastrantonisRubinstein2024}.
These approaches establish exponential-order lower bounds without the
sharp constant asserted in Theorem~\ref{thm:mahler}.

Exact results in substantial special classes also preceded the general
statement. Meyer and Reisner used shadow systems, one-parameter
deformations of convex bodies, to prove the sharp inequality and simplex
equality for polytopes with at most $n+3$ vertices. Their key input is
convexity of the reciprocal polar volume at the Santal\'o point along
such a deformation \cite[Theorems~1 and~10]{MeyerReisner2006}.
Kim and Reisner established strict
local minimality, with a quantitative estimate, at every simplex
\cite[Theorem~1]{KimReisner2011}; this is local information about the
space of convex bodies. In dimension three, Iriyeh and Shibata proved the
symmetric conjecture \cite{IriyehShibata2020}. The 2026 preprint of Chen,
Li, Xi, and Xu proves $P(K)\ge64/9$ for every three-dimensional
convex body, with equality exactly for tetrahedra
\cite{ChenLiXiXu2026}. Their shadow flows constrain the permitted
deformations so that they can treat arbitrary polytopes in dimension
three and then pass to general bodies. In contrast, the argument below
works with projections onto a fixed cone and their Gaussian averages.
Theorem~\ref{thm:mahler} concerns arbitrary convex bodies in all
dimensions, including its global equality classification.

\subsection*{Proof strategy}
The first step is the cone/Laplace correspondence developed by Klartag
\cite[Lemma~2.1 and Corollary~4.1]{Klartag2018}. Choose
$z_0\in\operatorname{int}K$ and put $m=n+1$. Lift the translated body
to the cone
\[
 C=\{(tx,t):t\ge0,\ x\in K-z_0\}\subset\mathbb R^m,
 \qquad D=\{y:\langle y,x\rangle\ge0\text{ for every }x\in C\}.
\]
For $V\in\operatorname{int}D$ and $U\in\operatorname{int}C$, set
$\chi_C(V)=\int_Ce^{-\langle V,x\rangle}\,dx$ and define $\chi_D(U)$
similarly. Section~\ref{sec:01-cones} shows that the desired inequality
follows from $\chi_C(V)\chi_D(U)\ge1$ whenever
$\langle U,V\rangle=m$. We use positive duality and retain the short
calculation to fix the signs and factorials.

The paired maps into $C$ and $D$ arise from Moreau's decomposition
\cite{Moreau1962,Soltan2019}. Let $\Pi_C$ and $\Pi_D$ be Euclidean nearest-point projections, and
write $Z=\Sigma^{1/2}G$, where $G=(G_1,\ldots,G_m)$ is standard Gaussian
and $\Sigma$ is positive definite. For a real layer parameter $z$,
choose a bias $\xi_z$ so that
$X_z=\Pi_C(Z+\xi_z)$ has mean $a(z)U$, where
$a(z)=\mathbb E\max\{g+z,0\}$ for a standard one-dimensional Gaussian $g$.
Its dual partner is $Y_z=\Pi_D(-Z-\xi_z)$.
Section~\ref{sec:02-projections} proves that every such bias exists and
then chooses linear coordinates and the covariance of $Z$
simultaneously. The dependence of the covariance on the projection
field is essential to the later matrix cancellation. Integrating the
projection maps over $z$ gives two injective, smoothed maps into the
cones. Change of variables and Jensen's inequality turn their Jacobians
into a lower bound for the normalized logarithm of the Laplace product
in Section~\ref{sec:03-entropy}. Write this lower bound as
$-\mathcal E$; the remaining task is to prove $\mathcal E\le0$.
Entropy and log-Laplace perturbations also appear in the functional
volume-product work of Fradelizi and Mar\'in Sola
\cite{FradeliziMarinSola2024}; that is a related analytical setting,
not an input theorem for these maps.

The main difficulty is that the derivative matrices $P_z=D\Pi_C(Z+\xi_z)$
need not commute and need not form an increasing family in $z$.
Estimating them independently would discard the information needed for
both the constant and equality. Sections~\ref{sec:04-quadratic}
and~\ref{sec:05-layers} therefore compare the full family to spectral
thresholds of a linear Gaussian matrix $L=\sum_{i=1}^mG_iM_i$, where the
$M_i$ are deterministic symmetric matrices. Classical Hermite expansion
and the Gaussian inverse-generator covariance identity
\cite[Proposition~2.1]{HoudrePrivault2002} separate this
linear part from higher Gaussian degrees. The resulting upper bound
for $\mathcal E$ subtracts a nonnegative layer defect and reduces
all other terms to functions of $L$ and the chosen covariance.

Section~\ref{sec:06-linear-equality} uses the matrix divided-difference
formula to express those functions as averages over pairs of eigenvalues.
A strict scalar inequality on each nontrivial interval absorbs the
remaining errors. Equality then forces the coefficient matrices $M_i$
to commute and the projection derivative to be diagonal in one fixed
basis. The cone splits into one-dimensional half-lines, so its bounded
section is a simplex. The converse is a direct volume computation.

The dimension-independent estimates used in this chain are stated in
Proposition~\ref{prop:scalar-input}. Appendix~\ref{sec:07-scalar} proves
the one-variable and pointwise layer bounds; Appendix~\ref{sec:08-segments}
proves the strict interval inequality. Both include the analytic
interpolation and tail arguments that turn their finite rational
certificates into statements on the whole real line. Numerical sampling
is not used as a substitute for these arguments. Figure~\ref{fig:proof-map}
records the logical dependencies and locates the independent scalar branch.

\begin{figure}[htbp]
\centering
\begin{tikzpicture}[
  font=\small,
  main/.style={draw=black!65,rounded corners=2pt,fill=black!3,
    align=center,text width=6.2cm,minimum height=1.12cm,inner sep=5pt},
  scalar/.style={draw=blue!45!black,rounded corners=2pt,fill=blue!4,
    align=center,text width=5.2cm,minimum height=1.3cm,inner sep=5pt},
  use/.style={-{Stealth[length=2.2mm]},draw=black!70,line width=.65pt}
]
\node[main] (cone) at (0,0)
  {Cone formulation\\Section~\ref{sec:01-cones}};
\node[main] (projection) at (0,-1.65)
  {Normalized Gaussian projections\\and entropy estimate\\
   Sections~\ref{sec:02-projections}--\ref{sec:03-entropy}};
\node[main] (quadratic) at (0,-3.3)
  {Quadratic identities and\\layer error bounds\\
   Sections~\ref{sec:04-quadratic}--\ref{sec:05-layers}};
\node[main] (linear) at (0,-4.95)
  {Linear-field estimate and\\equality characterization\\
   Section~\ref{sec:06-linear-equality}};
\node[align=center,font=\small\bfseries] at (9,.35)
  {Independent scalar input};
\node[scalar] (profiles) at (9,-1)
  {Scalar profile bounds\\
   Stated in Section~\ref{sec:01-cones};\\
   proved in Appendix~\ref{sec:07-scalar}};
\node[scalar] (segments) at (9,-4.95)
  {Segment variance estimate\\Appendix~\ref{sec:08-segments}};
\draw[use] (cone) -- (projection);
\draw[use] (projection) -- (quadratic);
\draw[use] (quadratic) -- (linear);
\draw[use] (profiles) -- (segments);
\draw[use] (profiles.west) -- (projection.east);
\draw[use] (profiles.south west) -| (5.2,-3.3) -- (quadratic.east);
\draw[use] (segments.west) -- (linear.east);
\end{tikzpicture}
\caption{Principal logical dependencies. Arrows indicate where inputs are
used, rather than the order of presentation. The scalar branch is independent
of the cone and dimension; its proofs follow the geometric argument.}
\label{fig:proof-map}
\end{figure}


\subsection*{Functional Mahler inequality}

The functional version replaces a convex body by a log-concave function
and polarity by convex conjugation. For a proper lower-semicontinuous
convex function $\varphi:\R^n\to\R\cup\{+\infty\}$, define its
Fenchel--Legendre transform by
\[
 \varphi^*(y)=\sup_{x\in\R^n}\{\langle x,y\rangle-\varphi(x)\},
 \qquad y\in\R^n.
\]
We use the convention $e^{-\infty}=0$.

\begin{corollary}[Functional Mahler inequality]\label{cor:functional-mahler}
For every integer $n\ge1$ and every proper lower-semicontinuous convex
function $\varphi:\R^n\to\R\cup\{+\infty\}$ such that
$0<\int_{\R^n}e^{-\varphi(x)}\,dx<\infty$,
\[
 \left(\int_{\R^n}e^{-\varphi(x)}\,dx\right)
 \left(\int_{\R^n}e^{-\varphi^*(y)}\,dy\right)\ge e^n.
\]
The second integral is allowed to be infinite, in which case the
inequality is immediate.
\end{corollary}

This sharp lower bound applies to general log-concave functions through
their convex potentials; no centering or normalization of $\varphi$ is
required. Its derivation uses Theorem~\ref{thm:mahler} in all dimensions,
not only in dimension $n$, through the implication of Fradelizi and
Meyer \cite[Proposition~2(a)]{FradeliziMeyer2008}.
Section~\ref{sec:consequences} gives the proof, an extremizing example,
and the distinction between geometric and functional equality cases.

The same functional bound has an entropy--transport consequence: for
full-dimensional log-concave probability measures whose densities vanish
at $\mathcal H^{n-1}$-almost every boundary point of their supports
(the essential-continuity condition), it bounds their summed entropies relative to
Lebesgue measure by a transport cost between the distributions of their
potential gradients, with additive constant $-3n$. The precise
definitions and statement appear in Corollary~\ref{cor:entropy-transport}.
This consequence uses the functional equivalence of Gozlan, as stated
in \cite[Theorem~1.3]{FradeliziGozlanZugmeyer2021}; essential continuity
is an additional hypothesis, not a requirement on the convex bodies in
Theorem~\ref{thm:mahler}.
